% SPDX-License-Identifier: GPL-3.0-or-later
% Copyright (C) 2026 Alessandro Greco
\chapter{Цензура, мосты и подключаемые транспорты}\label{cap:16-censura}
\aggiornatoal{23 сентября 2026 г.}{tor 0.4.9.13; pt-spec v1; obfs4proxy 0.0.14 (Debian); документация Tor Project}

\begin{obiettivi}
К концу этой главы вы сможете:
\begin{itemize}[nosep]
  \item описывать способы, которыми цензор может блокировать Tor;
  \item объяснять, что такое мосты и почему их распространение представляет сложную задачу;
  \item описывать протокол взаимодействия tor с подключаемым транспортом;
  \item различать стратегии obfs4, meek, Snowflake, WebTunnel и Conjure;
  \item эксплуатировать obfs4-мост в лабораторной сети и наблюдать, что способен видеть цензор (лабораторная работа~10).
\end{itemize}
\end{obiettivi}

Всё рассмотренное до сих пор предполагает, что клиент способен достичь сети. Во многих странах это не так: соединения с Tor блокируются, иногда только в периоды политической напряжённости, иногда постоянно. Эта глава завершает часть~II темой \textbf{обхода цензуры}, одним из наиболее активно развивавшихся направлений Tor Project в последние годы.

% ---------------------------------------------------------------------------
\section{Как блокируют Tor}\label{sec:cens-blocco}

Цензор, обычно государство, контролирующее провайдеров интернет-доступа страны, располагает несколькими инструментами:
\begin{description}
  \item[Блокировка адресов] Список реле публичен: это консенсус (\cref{cap:12-directory}). Цензор может просто загрузить его и заблокировать все IP-адреса реле либо только авторитетных серверов каталогов и реле, у которых клиенты получают консенсусы.
  \item[Блокировка сайтов] Блокировка \texttt{torproject.org} не позволяет людям загрузить Tor Browser, даже если сама сеть Tor остаётся доступной.
  \item[Распознавание протокола] \emph{Глубокий анализ пакетов} (deep packet inspection, DPI) позволяет распознать соединение Tor по особенностям рукопожатия TLS, размерам пакетов или временным характеристикам, даже когда адрес назначения ещё неизвестен.
  \item[Активное зондирование] Если DPI предполагает, что сервер может быть нераскрытым реле, цензор способен подключиться к нему и попытаться говорить по протоколу Tor. Если сервер отвечает, адрес блокируется.
  \item[Списки разрешённых адресов и протоколов] В крайнем случае цензор разрешает только явно одобренные протоколы или назначения и блокирует всё остальное.
\end{description}
Для каждой техники существует соответствующая контрмера, а противостояние цензоров и разработчиков продолжается. Такие проекты, как OONI, измеряют, где и как блокируется Tor (\cref{cap:24-metrics}).

% ---------------------------------------------------------------------------
\section{Мосты}\label{sec:cens-bridge}

Ответ на блокировку адресов --- \textbf{мост}: реле, отсутствующее в публичном консенсусе. Мост действует как входной узел (guard) (\cref{sec:guard}). Клиент использует его первым переходом, после чего цепь обычным образом продолжается по сети. Поскольку адрес моста не публикуется, цензор не может заблокировать его просто по списку. Мосты публикуют дескрипторы на \textbf{авторитетном сервере мостов (bridge authority)}, который проверяет достижимость и ведёт непубличный список. В опубликованной статистике мостов используется хеш \emph{от хеша} идентификатора, чтобы не раскрывать его.

Клиент настраивает мост с помощью \emph{строки конфигурации моста (bridge line)}. В лабораторной сети файл \file{torrc} клиента \texttt{test011bc} содержит:
\begin{lstlisting}[language=torrc,numbers=none]
UseBridges 1
Bridge obfs4 127.0.0.1:9909 2557828B596B33AB02075FE7EAD86F2DC3CEAFD5 cert=4sKDQih... iat-mode=0
ClientTransportPlugin obfs4 exec /usr/bin/obfs4proxy
\end{lstlisting}
Строка задаёт тип транспорта, адрес, отпечаток идентичности моста, используемый для проверки собеседника (\cref{cap:07-canali}), и параметры транспорта. Значение \texttt{cert=4sKDQih...} \emph{сокращено} для печатного примера и непригодно для непосредственного копирования: следует использовать полный сертификат и отпечаток, созданные собственной инстанцией моста.

\subsection{Проблема распространения}

Мост полезен, пока цензор не знает его адрес. Но для пользы этот адрес нужно сообщать пользователям, а цензор может притвориться пользователем. Идеального решения нет: распределитель должен предоставить достаточно мостов большому числу легитимных пользователей, раскрыв как можно меньше цензору. Tor Project распространяет мосты несколькими каналами, каждый со своей защитой от автоматизированного сбора~\cite{torsupport}:
\begin{itemize}
  \item веб-сайт \url{https://bridges.torproject.org/};
  \item электронная почта: \texttt{bridges@torproject.org};
  \item бот Telegram (\texttt{@GetBridgesBot});
  \item непосредственно через Tor Browser, который обнаруживает блокировку и запрашивает подходящие мосты для страны пользователя (\emph{Connection Assist}, на основе канала \emph{moat});
  \item частные каналы организаций, лично знающих своих пользователей.
\end{itemize}
Многие годы распределением мостов по каналам занимался BridgeDB. Его заменил \textbf{rdsys}; миграция завершилась в октябре 2024 года~\cite{torblogrdsys}. Операторы мостов могут выбрать канал опцией \opt{BridgeDistribution}.

\subsection{Lox}

\textbf{Lox}~\cite{tulloch2023lox} подходит к распространению через систему \emph{репутации}. Новый пользователь получает мост и анонимный криптографический сертификат (credential). Если мост со временем остаётся незаблокированным, уровень сертификата растёт вплоть до 4; начиная со второго уровня пользователь также может приглашать других. Если мост блокируется, можно получить новый, но репутация теряется. Идея состоит в том, что цензор, блокирующий каждый полученный мост, не способен повышать уровень, а честный пользователь способен. Такие удостоверения используют криптографические механизмы, которая не позволяет даже самой системе узнать, кто кого пригласил, защищая \emph{социальный граф} пользователей. Lox тестируется в альфа-версиях Tor Browser.

% ---------------------------------------------------------------------------
\section{Подключаемые транспорты}\label{sec:cens-pt}

Скрыть адрес недостаточно против распознавания протокола: сам трафик должен выглядеть иначе. Вместо изменения tor каждый раз, когда цензор учится распознавать новый признак, Tor Project определил интерфейс внешних программ, преобразующих трафик: \textbf{подключаемых транспортов} (pluggable transports, PT)~\spec{pt-spec}{Architecture Overview}.

\subsection{Архитектура}

PT --- отдельный процесс, который запускает tor, поэтому он называется \emph{управляемым (managed)}. На стороне клиента PT предоставляет tor локальный SOCKS-прокси. tor передаёт ему соединение к мосту, а PT преобразует трафик перед отправкой. На стороне моста PT слушает публичный порт, выполняет обратное преобразование и передаёт трафик локальному ORPort tor. Дополнительный небольшой протокол, \emph{Extended ORPort}, может передавать tor настоящий адрес клиента для статистики.

\begin{figure}[htbp]
\centering
% SPDX-License-Identifier: GPL-3.0-or-later
% Copyright (C) 2026 Alessandro Greco
%% Русское альтернативное описание EPUB.
\altfig{Упрощённая архитектура подключаемых транспортов: tor подключается к локальному SOCKS-прокси клиентского PT, который преобразует трафик. На мосту серверный PT обращает преобразование и передаёт соединение tor через ORPort либо, если он настроен, Extended ORPort. Промежуточные добровольческие прокси Snowflake не показаны. Цензор видит преобразованный трафик на контролируемом участке, но всё ещё может его распознать. Примеры: obfs4 напоминает случайные байты, Snowflake использует WebRTC, WebTunnel выглядит как HTTPS.}
% Архитектура pluggable transport (pt-spec)
\begin{tikzpicture}[font=\sffamily\scriptsize,
  proc/.style={draw,rounded corners=2pt,align=center,minimum height=0.8cm,minimum width=1.7cm},
  mac/.style={draw=torgrigio,dashed,rounded corners=4pt,inner sep=6pt}]
  \node[proc,fill=torviola!12] (tc) at (0,0) {tor\\(клиент)};
  \node[proc,fill=torverde!15] (pc) at (2.6,0) {клиентский PT};
  \node[proc,fill=torverde!15] (ps) at (8.2,0) {серверный PT};
  \node[proc,fill=torviola!12] (tb) at (10.8,0) {tor\\(bridge)};
  \node[mac,fit=(tc)(pc),label={[font=\sffamily\scriptsize]above:компьютер пользователя}] {};
  \node[mac,fit=(ps)(tb),label={[font=\sffamily\scriptsize]above:bridge}] {};
  \draw[->] (tc) -- node[above,font=\sffamily\tiny] {SOCKS} (pc);
  \draw[->] (ps) -- node[above,align=center,font=\sffamily\tiny] {ORPort /\\Extended ORPort} (tb);
  \draw[->,thick,decorate,decoration={snake,amplitude=1.2pt,segment length=6pt}] (pc) -- node[above] {преобразованный трафик} (ps);
  \node[server,minimum width=1.4cm,fill=avviso!10,draw=avviso] (cen) at (5.4,-1.5) {цензор};
  \draw[dotted,avviso] (cen) -- (5.4,-0.15);
  \node[text=torgrigio,align=center] at (5.4,-2.3) {байты без узнаваемой структуры (obfs4),\\WebRTC (Snowflake), HTTPS (WebTunnel)};
  \node[text=torgrigio,font=\sffamily\tiny,align=center,text width=10.5cm] at (5.4,-3.05) {Упрощённая схема: промежуточные добровольческие прокси Snowflake не показаны};
\end{tikzpicture}

\caption{Упрощённая архитектура подключаемого транспорта. На стороне клиента tor взаимодействует с PT через локальный SOCKS-прокси; на мосту PT передаёт восстановленный трафик tor через ORPort или, если он настроен, через Extended ORPort. Промежуточные добровольческие прокси Snowflake здесь опущены. На контролируемом участке цензор видит преобразованный трафик, однако возможность его распознавания не исключена.}
\label{fig:pt}
\end{figure}

\subsection{Протокол между tor и PT}

tor настраивает PT переменными окружения, а PT отвечает строками в стандартный вывод~\spec{pt-spec}{Specification}. Это можно проверить, вручную запустив \cmd{obfs4proxy} в лабораторном контейнере точно так, как это сделал бы tor. В клиентском режиме:
\begin{lstlisting}[language=bash,numbers=none]
$ env TOR_PT_MANAGED_TRANSPORT_VER=1 TOR_PT_STATE_LOCATION=/tmp/pt \
      TOR_PT_EXIT_ON_STDIN_CLOSE=1 TOR_PT_CLIENT_TRANSPORTS=obfs4 obfs4proxy
VERSION 1
CMETHOD obfs4 socks5 127.0.0.1:40057
CMETHODS DONE
\end{lstlisting}
PT объявляет версию протокола и сообщает, что предоставляет транспорт \texttt{obfs4} как прокси SOCKS5 на выбранном им локальном порту. В серверном режиме с переменными, задающими публичный адрес прослушивания и адрес назначения ORPort:
\begin{lstlisting}[language=bash,numbers=none]
$ env ... TOR_PT_SERVER_TRANSPORTS=obfs4 \
      TOR_PT_SERVER_BINDADDR=obfs4-127.0.0.1:19999 \
      TOR_PT_ORPORT=127.0.0.1:5104 obfs4proxy
VERSION 1
SMETHOD obfs4 127.0.0.1:19999 ARGS:cert=5W6vcbWbFGGZgabwm4C3WipezwnAfadoBhqfNVpY72xnkjB/Szczbxg8fLWJJ1uRlyjPSg,iat-mode=0
SMETHODS DONE
\end{lstlisting}
Параметры после \texttt{ARGS:} нужны клиенту и включаются в строку конфигурации моста. PT также записывает их в \file{obfs4\_bridgeline.txt} своего каталога состояния. Переменная \file{TOR\_PT\_EXIT\_ON\_STDIN\_CLOSE} велит PT завершиться, когда tor закрывает входной поток, чтобы транспорт не пережил запустивший его процесс. (Из-за недосмотра пример спецификации утверждает, что PT пишет сообщения в \enquote{stdin}; на практике, как мы только что увидели, он пишет их в стандартный вывод.)

Команда сервера \texttt{env ...} приведена \emph{схематически}: многоточие обозначает другие необходимые переменные окружения и не должно вводиться буквально. Порт SOCKS5 в ответе клиентского транспорта также взят из эталонного запуска и не является постоянным.

\subsection{Lyrebird}

Транспортная программа, используемая Tor Browser, --- \textbf{lyrebird}, происходящая от более раннего \cmd{obfs4proxy}. Она реализует obfs4, meek, Snowflake и WebTunnel~\cite{torsupport}. В лаборатории используется \cmd{obfs4proxy} 0.0.14 из пакетов Debian, который по-прежнему совместим с obfs4.

% ---------------------------------------------------------------------------
\section{Стратегии}\label{sec:cens-strategie}

Современные транспорты следуют двум противоположным стратегиям: \emph{не быть похожими ни на что распознаваемое} либо \emph{быть похожими на что-то разрешённое}.

\subsection{obfs4: не быть похожим ни на что распознаваемое}

\textbf{obfs4}, разработанный Yawning Angel, превращает трафик в последовательность байтов, неотличимых от случайных данных: без узнаваемого заголовка, без рукопожатия TLS и с переменными размерами пакетов. Его эффективность основана на трёх решениях:
\begin{itemize}
  \item Обмен ключами представляет собой вариант ntor (\cref{sec:ntor}), где открытые ключи Curve25519 кодируются с помощью Elligator2, техники, делающей точку кривой неотличимой от случайной строки.
  \item Клиент должен знать параметр \texttt{cert} моста, содержащий идентификатор и открытый ключ. Без него наблюдатель не может построить корректное рукопожатие, а мост не отвечает. Это противодействует \textbf{активному зондированию}: любопытный цензор, подключившийся к мосту, не получает распознаваемого ответа.
  \item Параметр \texttt{iat-mode} (\emph{inter-arrival time}) ценой производительности может добавлять случайные задержки и фрагментацию, дополнительно скрывая временные характеристики.
\end{itemize}
Ограничение obfs4 является одновременно его сильной стороной: цензор со списком разрешённых протоколов или блокировкой \enquote{всего, что выглядит случайным} способен его остановить. Так происходит в некоторых странах.

\subsection{meek: за крупным сервисом}

\textbf{meek}~\cite{fifield2015fronting} использует \emph{маскировку домена (domain fronting)}. Клиент открывает HTTPS-соединение к крупной сети доставки контента, указывая безобидное доменное имя в видимой части соединения, а настоящий пункт назначения, сервер-шлюз в Tor, сообщает только внутри зашифрованных данных. Блокировка meek потребовала бы заблокировать весь сервис с большим сопутствующим ущербом. Однако техника зависит от разрешения провайдеров, и многие в последние годы сделали её невозможной. Кроме того, meek медленный и дорогой, поскольку весь трафик проходит через серверы провайдера.

\subsection{Snowflake: множество эфемерных прокси}

\textbf{Snowflake}~\cite{bocovich2024snowflake} использует тысячи добровольцев, запускающих небольшие прокси, например расширение браузера или программу на сервере. Клиент связывается с \emph{посредником (broker)} через трудно блокируемый канал, например domain fronting, получает адрес прокси и подключается к нему через WebRTC, протокол видеозвонков в браузерах. Прокси пересылает трафик на мост Snowflake Tor Project. Прокси много, и они постоянно меняются. Заблокировать все их адреса непрактично, а блокировка WebRTC означает блокировку видеозвонков.

\subsection{WebTunnel: выглядеть как веб-сайт}

\textbf{WebTunnel}, объявленный Tor Project в марте 2024 года, заворачивает трафик в HTTPS-соединение с поведением, похожим на WebSocket~\cite{torblogwebtunnel}. Мост может находиться за обычным сайтом, деля с ним домен, адрес и порт. Обратный прокси передаёт мосту только запросы к секретному пути, а всем остальным обслуживает обычный сайт. WebTunnel рассчитан на цензоров со списками разрешённых протоколов, против которых obfs4 не работает: он имеет \enquote{форму} разрешённого протокола.

\subsection{Conjure: скрытие в адресном пространстве}

\textbf{Conjure}~\cite{frolov2019conjure} --- система \emph{преломление сетевого трафика (refraction networking)}. Вместо моста по известному адресу она использует станции, установленные у сотрудничающих провайдеров интернет-доступа. Эти станции наблюдают трафик к диапазонам неиспользуемых адресов и перенаправляют его в Tor. Клиент подключается к \enquote{фантомному} адресу, выбранному из миллионов вариантов. Чтобы заблокировать Conjure, цензору пришлось бы блокировать целые части Интернета. По документации Tor Project, Conjure успешно тестируется, но всё ещё требует стабилизации и доступен только в альфа-версиях Tor Browser и некоторых приложениях~\cite{torsupport}.

\begin{attenzione}[Ни один транспорт не работает везде]
У каждого транспорта есть слабости, а его эффективность зависит от страны и меняется со временем. Например, на форуме Tor Project в апреле 2025 года пользователи сообщали, что obfs4, meek и Snowflake не работали в Китае, а WebTunnel подключался, но быстро блокировался~\cite{forumcina2025}. Поэтому Tor Browser предлагает несколько транспортов и функцию выбора подходящего для страны пользователя.
\end{attenzione}

% ---------------------------------------------------------------------------
\section{Лабораторная работа: obfs4-мост глазами цензора}\label{sec:lab-censura}

\begin{laboratorio}[10]{Obfs4-мост в лабораторной сети}
\textbf{Цель.} Запустить сеть с obfs4-мостом и использующим его клиентом; захватить и сравнить трафик обычного Tor-соединения и obfs4-соединения.\\
\textbf{Предварительные условия.} Лабораторная сеть (\cref{cap:19-rete-privata}). Образ содержит \cmd{obfs4proxy} и \cmd{snowflake-client} из пакетов Debian.\\
\textbf{Оценочное время.} 1--2 часа.
\end{laboratorio}

\paragraph{Шаг 1: сеть с мостом.} Файл \file{labs/ru/лаборатория/chutney/networks/book-bridge} добавляет к сети книги авторитетный сервер мостов (\texttt{test003ba}), obfs4-мост (\texttt{test009br}) и использующий его клиент (\texttt{test011bc}):
\begin{lstlisting}[language=bash,numbers=none]
$ make -C labs/ru/лаборатория clean
$ make -C labs/ru/лаборатория start NETWORK=book-bridge
\end{lstlisting}
Журнал клиента из предыдущего эталонного запуска содержит:
\begin{lstlisting}[numbers=none]
[notice] new bridge descriptor 'test009br' (fresh): $2557828B...~test009br
[warn] Only one bridge (transport: 'obfs4') is configured. You should
       have at least two for conflux, for any transport that is not 'snowflake'.
[notice] Bootstrapped 100% (done): Done
\end{lstlisting}
Клиент получил дескриптор моста через сам мост и завершил загрузку конфигурации. Предупреждение напоминает, что Conflux (\cref{cap:31-futuro}) использует две параллельные цепи и не может делать это с единственным мостом. Веб-запрос через SOCKS-порт этого клиента (9011) работает так же, как через обычный клиент.

\paragraph{Шаг 2: что видит цензор.} Внутри контейнера захватываем трафик \cmd{tcpdump} на obfs4-порту моста (9909) и обычном ORPort (5104), пока клиент использует мост и мы открываем канал к \texttt{test004r} лабораторной работой~1. Затем анализируем захваченные пакеты с помощью \cmd{tshark}. Вот результат одного из наших сеансов:
\begin{lstlisting}[numbers=none]
== ORPort 5104 (TLS), первые байты от клиента ==
1533  16030105f8010005f403033e1234a0401520d1c52b...
== obfs4 9909, первые байты от клиента ==
1957  b566475349099523c53005e1004a392160f15d473e...
== пакеты, распознанные как TLS ==
порт 5104: 83
порт 9909: 0
== размеры сегментов obfs4 (клиент -> мост) ==
1957 1677 1888 1891 1716 1716 3436 1395 1716 1748 1500 1194 ...
\end{lstlisting}
На ORPort первые байты однозначны: \texttt{16 03 01} начинает запись рукопожатия TLS, а Wireshark распознаёт весь разговор как TLS. Цензор видит TLS-соединение, которое можно исследовать на признаки Tor; адрес реле также опубликован в консенсусе. На obfs4-мост байты не имеют узнаваемой структуры, ни один пакет не определяется как TLS, а размеры сегментов не повторяют 514-байтовые ячейки. Цензор видит лишь поток кажущихся случайными данных к неизвестному адресу.

\paragraph{Шаг 3: зондирование моста.} Пытаемся подключиться к порту 9909 моста через лабораторную работу~1, как мог бы сделать цензор, подозревающий реле:
\begin{lstlisting}[numbers=none]
TimeoutError: The handshake operation timed out
\end{lstlisting}
Мост не отвечает \emph{вообще}: ни сообщения об ошибке, ни одного байта, пока клиент не сдаётся через 15 секунд. Для obfs4 наш TLS ClientHello --- некорректное рукопожатие, а реакция на некорректное рукопожатие --- молчание. Цензор, зондирующий адрес, не узнаёт ничего, даже того, что там слушает сервис иного протокола. Без параметра \texttt{cert} он всё равно не способен сформировать корректное рукопожатие obfs4.

% ---------------------------------------------------------------------------
\section*{Итоги}
\begin{itemize}
  \item Цензор может блокировать Tor по адресам, распознавать протокол, зондировать подозрительные серверы или разрешать только явно одобренный трафик.
  \item Мосты --- непубличные реле. Их распространение, сегодня управляемое rdsys и исследуемое посредством механизма репутации Lox, балансирует доступ пользователей и обнаружение цензорами.
  \item Подключаемые транспорты --- внешние процессы, преобразующие трафик; tor настраивает их переменными окружения и читает ответы (\texttt{CMETHOD}, \texttt{SMETHOD}).
  \item obfs4 делает трафик похожим на случайный и противостоит активному зондированию; meek, Snowflake, WebTunnel и Conjure скрывают его за сервисами, протоколами или адресами, блокировка которых дорого обходится цензору.
  \item Ни один транспорт не работает везде и навсегда.
\end{itemize}

% ---------------------------------------------------------------------------
\section*{Упражнения}

\begin{esercizio}[1]\label{es:16-guard}
Чем мост отличается от входного узла? Почему клиент, использующий мосты, не использует дополнительно публичный входной узел?
\end{esercizio}

\begin{esercizio}[1]\label{es:16-strategie}
Классифицируйте obfs4, meek, Snowflake, WebTunnel и Conjure по стратегии (\enquote{не быть похожим ни на что распознаваемое} или \enquote{быть похожим на что-то разрешённое}) и укажите, что пришлось бы блокировать цензору, чтобы остановить каждый из них.
\end{esercizio}

\begin{esercizio}[2]\label{es:16-sondaggio}
Опишите, как цензор может обнаружить obfs4-мост, наблюдая только трафик и не прибегая к активному зондированию. Какие статистические признаки он может использовать и как \texttt{iat-mode} пытается им противодействовать?
\end{esercizio}

\begin{esercizio}[2]\label{es:16-lox}
Почему в Lox цензор не способен повышать уровень репутации? Что он мог бы сделать, чтобы обойти механизм, и почему это дорого?
\end{esercizio}

\begin{esercizio}[2]\label{es:16-entropia}
Вычислите энтропию Шеннона в битах на байт для первых 10\,KB, захваченных из obfs4-соединения и обычного TLS-соединения в лаборатории. Что вы ожидаете и почему разница может оказаться небольшой? Что же тогда делает TLS распознаваемым?
\end{esercizio}

\begin{esercizio}[3]\label{es:16-sondaggio-lab}
Попробуйте запустить лабораторную работу~1 против obfs4-порта своего моста. Какая ошибка возникает? Затем используйте \cmd{obfs4proxy} как клиент с неправильным параметром \texttt{cert}: что видно в журнале моста?
\end{esercizio}

% ---------------------------------------------------------------------------
\section*{Дополнительная литература}
\begin{itemize}
  \item \titolo{Pluggable Transport Specification (Version 1)} (pt-spec)~\cite{torspec}.
  \item L.~Tulloch и I.~Goldberg, \titolo{Lox: Protecting the Social Graph in Bridge Distribution}~\cite{tulloch2023lox}.
  \item D.~Fifield et al., \titolo{Blocking-resistant communication through domain fronting}~\cite{fifield2015fronting}.
  \item C.~Bocovich et al., \titolo{Snowflake, a censorship circumvention system using temporary WebRTC proxies}~\cite{bocovich2024snowflake}.
  \item S.~Frolov et al., \titolo{Conjure: Summoning Proxies from Unused Address Space}~\cite{frolov2019conjure}.
  \item Запись блога Tor Project о запуске WebTunnel~\cite{torblogwebtunnel} и страницы поддержки мостов~\cite{torsupport}.
\end{itemize}
